\section{Evidence and Provenance Boundary}
\label{app:evidence-provenance}

The maintained evidence is retained under:
\begin{itemize}
 \item \path{calculations/ICMSEP2026/conference/paper_1/isolated-band/};
 \item \path{calculations/ICMSEP2026/}\\
       \path{conference/paper_1/multiband-alignment/};
 \item \path{calculations/research-monograph/periodic-2d/}; and
 \item \path{calculations/ICMSEP2026/conference/paper_1/}.
\end{itemize}
Each directory contains its protocol, result records, independent verification
code, figures, report, and checksum manifest. The evidence has been classified
as controlled illustrative software and numerical verification. It is not a
public data deposit, DOI, material validation, or uncertainty analysis.

The isolated-band package contains the prospectively frozen M1 lowest-band
calculation. Its historical identity does not substitute for a retained gap or
band-isolation result. The multiband package contains the prospectively frozen
M2 rank-two alignment and
locality calculation, including its independent reconstruction and disjoint
withheld diagnostics. The two-dimensional reports contain additional composite
and topology controls.
The conference-paper parent calculation contains the direct
quadratic admissible-set map, exact common witness, rational separation
certificate, independent verifier, and figure inputs. Only the bounded results
directly used in the paper are summarized in its main text. Historical external
executions are not rerun by the manuscript.

\subsection{Verifier and digest semantics}

The standalone M1 and M2 retained verifiers decode the frozen JSON documents
and reconstruct their finite protocols without importing the corresponding
producer Actions. They remain dependent on shared numerical libraries,
runtime behavior, and declared mathematical conventions; passing establishes
bounded reconstruction consistency rather than an independent physical oracle.
The checksum manifests bind retained bytes to recorded SHA-256 digests and
therefore detect inconsistency relative to those manifests. A digest manifest
alone is not an external trusted timestamp and does not, without an independent
commit or archival record, prove when the bytes were created.

\subsection{Curation work}

The remaining evidence-curation tasks are:
\begin{itemize}
 \item[$\square$] trace every numerical statement and plotted value in the
 paper to a retained result field or report table;
 \item[$\square$] verify the figure sources against the retained checksum
 manifests;
 \item[$\square$] record the exact software revision and runtime environment
 associated with each retained calculation;
 \item[$\square$] decide which machine-readable records and figure-generation
 inputs belong in a public archival package; and
 \item[$\square$] obtain authorization before publication or deposition.
\end{itemize}
These practices follow general reproducible-computation guidance~\cite{sandve2013}.
Completing the tasks strengthens provenance and publication readiness; it does
not change the evidence class from controlled numerical verification to
material validation.
